% gerado por ferramentas/md2tex_enunciado.py a partir de questoes/q128.md (não editar à mão)
Ao trafegar à noite, é necessário manter os faróis dos veículos acesos. Quando a umidade relativa do ar é baixa, o motorista consegue ter boa visibilidade da pista. No entanto, em noites em que há neblina, essa visibilidade é reduzida.\par
Nessa situação, a visibilidade reduzida deve-se à\par
\begin{alternativas}
  \item dispersão da luz pelas gotículas de água suspensas no ar.
  \item polarização da luz pelas gotículas de água dispersas no ar.
  \item condensação das gotículas de água na superfície do farol.
  \item solubilização das gotículas de água presentes no ar.
  \item redução da potência da luz emitida pelo farol.
\end{alternativas}
